% !PRINTED 300
\weberpage{300}
%<<CONT
und ihre Potenzen; ebenso die Substitutionen
\[
\begin{pmatrix} 1, & 2 \\ 0, & 1 \end{pmatrix},\quad \begin{pmatrix} 2, & 1 \\ -1, & 0 \end{pmatrix},
\]
und folglich auch
\weq{6}{B = \begin{pmatrix} 0, & 1 \\ -1, & 0 \end{pmatrix} = \begin{pmatrix} 1, & 2 \\ 0, & 1 \end{pmatrix}\begin{pmatrix} 2, & 1 \\ -1, & 0 \end{pmatrix}.}
Wenn aber die Substitutionen $A$, $B$ in $\mathfrak{K}$ enthalten sind, so muß $\mathfrak{K}$ nach § 79 mit $\mathfrak{L}_0$ identisch sein. Wir haben also den Satz:

\gesperrt{Der Index eines Teilers von $\mathfrak{L}_0$ kann nicht kleiner als $p$ sein.}

Unser Problem beschränkt sich also auf die Frage nach den Teilern von $\mathfrak{L}_0$ vom Index $p$ oder vom Grade $\dfrac{p^2-1}{2}$.

Es sei $\mathfrak{K}$ ein solcher Teiler und $S_0, S_1, \ldots, S_{\nu-1}$ seien seine Elemente, so daß die Zahl $\nu$ den Wert
\weq{7}{\nu = \frac{p^2-1}{2}}
hat.

Da der Grad eines jeden Elementes einer Gruppe immer ein Teiler des Grades der Gruppe ist, so kann, da $\nu$ durch $p$ nicht teilbar ist, in $\mathfrak{K}$ \gesperrt{kein} Element vom Grade $p$ vorkommen; nehmen wir also für $T$ die Substitution vom Grade $p$:
\weq{8}{A = \begin{pmatrix} 1, & 0 \\ 1, & 1 \end{pmatrix},}
so läßt sich die Gruppe $\mathfrak{L}_0$ in die $p$ Nebengruppen zerlegen:
\weq{9}{\mathfrak{L}_0 = \mathfrak{K},\ A\mathfrak{K},\ A^2\mathfrak{K}\ \ldots,\ A^{p-1}\mathfrak{K}.}
Mit $\mathfrak{K}$ sind von demselben Grade, also auch von demselben Index alle seine konjugierten Teiler
\[
T\mathfrak{K}T^{-1}.
\]
Es sei nun $g$ eine primitive Wurzel der Primzahl $p$, und
\weq{10}{C = \begin{pmatrix} g, & 0 \\ 0, & g^{-1} \end{pmatrix},}
eine Substitution, die offenbar vom Grade $\dfrac{p-1}{2}$ ist. Diese muß, wie jede Substitution von $\mathfrak{L}_0$, in einer der Reihen (9) vorkommen; d.\,h. es gibt eine Substitution $S$ in $\mathfrak{K}$ und einen Exponenten $\lambda$, für den
\weq{11}{C = A^\lambda S,\quad S = A^{-\lambda} C.}
